\documentclass[10pt]{article}
\usepackage[T1]{fontenc}
\usepackage[utf8]{inputenc}
\usepackage[french]{babel}
\usepackage{amsmath,amssymb,mathrsfs,mathpazo}
\usepackage{geometry}
\geometry{a4paper,margin=2cm}
\usepackage{graphicx,array,tabularx,tabularray,colortbl,multirow,multicol,booktabs,enumitem,pifont,ulem,fancybox,diagbox,cancel}
\usepackage{pstricks,pstricks-add,pst-plot,pst-tree,pst-node,pst-text,pst-func,pst-math,pst-bspline,pst-3dplot,pst-eucl,pst-solides3d}
\let\pscalculate\undefined
\usepackage{pst-bezier}
\usepackage{tikz,pgfplots,tkz-tab}
\usetikzlibrary{arrows,arrows.meta,patterns,calc,positioning,trees,decorations.markings,shapes.geometric,shapes.misc,shapes.arrows,intersections,decorations.pathmorphing,decorations.pathreplacing,backgrounds,shadows,babel}
\pgfplotsset{compat=1.18}
\usepgfplotslibrary{fillbetween}
\usepackage[np]{numprint}
\usepackage{listings,makecell,pas-tableur,twoopt}
\let\barre\undefined
\usepackage[pstricks]{tablvar}
\definecolor{darkgreen}{RGB}{0,100,0}
\providecommand{\red}{\color{red}}
\providecommand{\blue}{\color{blue}}
\providecommand{\green}{\color{green}}
\providecommand{\black}{\color{black}}
\providecommand{\white}{\color{white}}
\providecommand{\gray}{\color{gray}}
\providecommand{\cyan}{\color{cyan}}
\providecommand{\magenta}{\color{magenta}}
\providecommand{\yellow}{\color{yellow}}
\providecommand{\violet}{\color{violet}}
\providecommand{\orange}{\color{orange}}
\providecommand{\euro}{\texteuro}
\providecommand{\EUR}{\texteuro}
\providecommand{\ieme}{\textsuperscript{e}}
\providecommand{\degres}{\ensuremath{^\circ}}
\providecommand{\degree}{\ensuremath{^\circ}}
\providecommand{\vv}[1]{\overrightarrow{#1}}
\providecommand{\Coord}[1]{#1}
\providecommand{\decosix}{}
\providecommand{\starredbullet}{\ensuremath{\bullet}}
\pagestyle{empty}
\setlength{\parindent}{0pt}
\usepackage{adjustbox}\DeclareUnicodeCharacter{2236}{:}\DeclareUnicodeCharacter{2190}{\ensuremath{\leftarrow}}
\providecommand{\R}{\mathbb{R}}
\providecommand{\N}{\mathbb{N}}
\providecommand{\Z}{\mathbb{Z}}
\providecommand{\Q}{\mathbb{Q}}
\providecommand{\D}{\mathbb{D}}
\providecommand{\C}{\mathbb{C}}
\providecommand{\e}{{\mathrm{e}}}
\providecommand{\vect}[1]{\overrightarrow{#1}}
\providecommand{\vectt}[1]{\overrightarrow{\mathrm{#1}}}
\providecommand{\vvt}[1]{\overrightarrow{\mathrm{#1}}}
\providecommand{\np}[2][]{{#2~#1}}
\providecommand{\vv}[1]{{\overrightarrow{#1}}}
\renewcommand{\d}{{\mathrm{d}}}
\providecommand{\degres}{{^\circ}}
\providecommand{\degre}{{^\circ}}
\providecommand{\blue}{\color{blue}}
\providecommand{\red}{\color{red}}
\providecommand{\enskip}{\quad}
\providecommand{\llbracket}{\lbrack\!\lbrack}
\providecommand{\rrbracket}{\rbrack\!\rbrack}
\providecommand{\euro}{{\text{€}}}
\providecommand{\no}{{\text{n°}}}
\providecommand{\interval}[2]{{[#1;#2]}}
\providecommand{\intervFO}[2]{{[#1;#2[}}
\providecommand{\dotfill}{\ldots}

\providecommand{\ds}{\displaystyle}
\providecommand{\textding}[1]{#1}
\providecommand{\Oij}{\(\mathrm{O};\vec{\imath},\vec{\jmath}\)}
\providecommand{\Oijk}{\(\mathrm{O};\vec{\imath},\vec{\jmath},\vec{k}\)}
\providecommand{\Ouv}{\(\mathrm{O};\vec{u},\vec{v}\)}
\providecommand{\intervFF}[2]{\ensuremath{[#1;#2]}}
\renewcommand{\euro}{\ifmmode\text{\texteuro}\else\texteuro\fi}
\DeclareUnicodeCharacter{00B0}{\ifmmode^{\circ}\else\textdegree\fi}
\DeclareUnicodeCharacter{20AC}{\ifmmode\text{\texteuro}\else\texteuro\fi}
% Énoncés : Éducation nationale. Transcription et corrigés : APMEP, Denis Vergès et contributeurs. Découpage et classement : CoopMaths, Sébastien Lozano, Éric Elter et contributeurs. Adaptation de la présentation : Maths Explorer.
% CC BY-SA : https://coopmaths.fr/www/about/
% https://coopmaths.fr/alea/static/bac/2023/tex/bac_2023_03_sujet2_etrangers_groupe2_1.tex
\begin{document}
\noindent\textit{APMEP / CoopMaths. CC BY-SA. Présentation adaptée par Maths Explorer.}\par\medskip
\medskip

On considère la fonction $f$ définie sur $\R$ par:
\[f(x) = \dfrac{1}{1 + \text{e}^{- 3x}}.\]

On note $\mathcal{C}_f$ sa courbe représentative dans un repère orthogonal du plan.

On nomme A le point de coordonnées $\left(0~;~\dfrac12\right)$ et B le point de coordonnées $\left(1~;~\dfrac54\right)$.

On a tracé ci-dessous la courbe $\mathcal{C}_f$ et $\mathcal{T}$ la tangente à la courbe $\mathcal{C}_f$ au point d'abscisse $0$.

\begin{center}
\psset{xunit=1.25cm,yunit=2.5cm,arrowsize=2pt 3}
\begin{pspicture*}(-2.5,-0.5)(3,1.5)
\psgrid[subgriddiv=4](-3,-1)(3,2)
\psaxes[linewidth=1.25pt,Dx=10,Dy=10]{->}(0,0)(-2.5,-0.5)(3,1.5)
\psaxes[linewidth=1.25pt,Dx=10,Dy=10](0,0)(1.1,1.1)
\psplot[plotpoints=2000,linewidth=1.25pt,linecolor=blue]{-2.5}{3}{2.71828 x 3 mul exp 2.71828 x 3 mul exp 1 add div}
\psplot[plotpoints=2000,linewidth=1.25pt]{-2.5}{3}{3 x mul 4 div 0.5 add}
\uput[l](0,0.5){A}\uput[dr](1,1.25){B}
\uput[dr](1,0){1}\uput[ul](0,1){1}
\psdots[dotstyle=+,dotangle=45,dotscale=1.5](0,0.5)(1,1.25)
\uput[d](2.8,0.9){\blue $\mathcal{C}_f$}\uput[r](-1.2,-0.4){$\mathcal{T}$}
\end{pspicture*}
\end{center}

\textbf{Partie A : lectures graphiques}

\medskip

Dans cette partie, les résultats seront obtenus par lecture graphique. Aucune justification n'est demandée.

\medskip

\begin{enumerate}
\item Déterminer l'équation réduite de la tangente $\mathcal{T}$.
\item Donner les intervalles sur lesquels la fonction $f$ semble convexe ou concave.
\end{enumerate}

\bigskip

\textbf{Partie B : étude de la fonction}

\medskip

\begin{enumerate}
\item On admet que la fonction $f$ est dérivable sur $\R$.

Déterminer l'expression de sa fonction dérivée $f'$.
\item Justifier que la fonction $f$ est strictement croissante sur $\R$.
\item
	\begin{enumerate}
		\item Déterminer la limite en $+ \infty$ de la fonction $f$.
		\item Déterminer la limite en $- \infty$ de la fonction $f$.
	\end{enumerate}
\item Déterminer la valeur exacte de la solution $\alpha$ de l'équation $f(x) = 0,99$.
\end{enumerate}

\bigskip

\textbf{Partie C : Tangente et convexité}

\medskip

\begin{enumerate}
\item Déterminer par le calcul une équation de la tangente $\mathcal{T}$ à la courbe $\mathcal{C}_f$ au point d'abscisse $0$.

On admet que la fonction $f$ est deux fois dérivable sur $\R$. 

On note $f''$ la fonction dérivée seconde de la fonction $f$.

On admet que $f''$ est définie sur $\R$ par :
\[f''(x) = \dfrac{9\text{e}^{-3x} \left(\text{e}^{-3x} - 1\right)}{\left(1 + \text{e}^{-3x} \right)^3}.\]

\item Étudier le signe de la fonction $f''$ sur $\R$.
\item 
	\begin{enumerate}
		\item Indiquer, en justifiant, sur quel(s) intervalle(s) la fonction $f$ est convexe.
		\item Que représente le point A pour la courbe $\mathcal{C}_f$ ?
\item  En déduire la position relative de la tangente $\mathcal{T}$ et de la courbe $\mathcal{C}_f$.

Justifier la réponse.
	\end{enumerate}
\end{enumerate}

\medskip

\end{document}
